\documentclass[10pt,a4paper]{article}
\usepackage[top=0.55in,bottom=0.55in,left=0.6in,right=0.6in]{geometry}
\usepackage[T1]{fontenc}
\usepackage{times,microtype,enumitem,hyperref,titlesec}
\pagestyle{empty}
\setlength{\parindent}{0pt}
\setlength{\parskip}{0pt}
\linespread{1.04}
\hypersetup{hidelinks,breaklinks=true}
\titleformat{\section}{\bfseries\normalsize}{}{0pt}{}[\titlerule]
\titlespacing*{\section}{0pt}{10pt}{4pt}
\newcommand{\entry}[3]{\noindent\textbf{#1}\hfill\textbf{#3}\\[-1pt]\textit{#2}\par}
\newenvironment{cvitems}{\begin{itemize}[leftmargin=1.1em,label=\textbullet,itemsep=2pt,topsep=3pt,parsep=0pt,partopsep=0pt]}{\end{itemize}}
\begin{document}
{\LARGE\bfseries Aarush Ram Anandh}\par\vspace{4pt}
\href{mailto:aarushramanandh@gmail.com}{aarushramanandh@gmail.com}\quad|\quad +91 9659786783\par\vspace{2pt}
\href{https://github.com/MushuRam01}{github.com/MushuRam01}\quad|\quad
\href{https://www.linkedin.com/in/aarush-ram-anandh/}{linkedin.com/in/aarush-ram-anandh}\par

\section{EXECUTIVE SUMMARY}
Machine learning researcher working on diffusion models and controllable optimal transport.\\
Public preprint, C++/Lean research implementation, and controlled reinforcement learning studies.\\
Python/PyTorch research across vision, RL and asthma detection at approximately 70\% sensitivity.\par
\section{EDUCATION}
\textbf{Ashoka University} --- B.Sc. (Hons.), Computer Science and Mathematics \hfill \textbf{Expected 2028}
\section{RESEARCH EXPERIENCE}
\entry{Constrained Image-Synthesis Lab, Ashoka University}{Research Fellow}{Sep 2025--Present}
\begin{cvitems}
  \item Improved latent-diffusion colour fidelity to $\Delta E_{00}<1$, with approximately \textbf{4$\times$ better colour accuracy} than evaluated open-source baselines; results reported in the preprint below.
  \item Developing \textbf{Controlled Generalised Sliced Wasserstein (CGSW)}, an optimal-transport framework with control-conditioned probes, projections and slice aggregation; paper submitted to \textbf{JMLR}.
  \item Implementing \textbf{C++} methods and machine-verified \textbf{Lean} specifications/proofs for CGSW.
  \item Conducting a \textbf{CGSW experiment} on hyperspectral waste images with \textbf{224 spectral bands and seven categories}, constructing object-level spectral measures to study class separability.
\end{cvitems}
\entry{Koita Centre for Digital Health, Ashoka University}{Research Intern (part-time)}{May 2026--Present}
\begin{cvitems}
  \item Contributed to a wearable-data pipeline achieving approximately \textbf{70\% sensitivity (recall)} for \textbf{asthma exacerbation detection} in study evaluation.
  \item Built HR/HRV smoothing, residual and change-point features for \textbf{191 patients}; implemented patient-wise 80:20 splits, five seeds and patient-bootstrap confidence intervals.
\end{cvitems}
\section{SELECTED PROJECTS}
\entry{Reinforcement Learning for Pokemon Battles}{Two-Person RL Coursework Project \href{https://github.com/Kabir1607/RL_MidSem_Kabir_Aarush}{[GitHub]}}{Ongoing}
\begin{cvitems}
  \item Co-developed a \textbf{Gymnasium/poke-env} environment with \textbf{763 observation features}, legality masks over 26 actions, potential-based reward shaping, and \textbf{MaskablePPO and masked DQN} agents.
  \item Evaluated PPO across \textbf{two variants and five seeds each} at 200,000 steps/run; achieved \textbf{99.7--100\% wins} against the fixed heuristic (1,020 final battles/run), with convergence and uncertainty analysis.
\end{cvitems}
\entry{MEIKURAL: Hindustani Raag Recognition}{Audio Classification / Signal Processing \href{https://github.com/MushuRam01/MEIKURAL}{[GitHub]}}{Ongoing}
\begin{cvitems}
  \item Developing a \textbf{Python/PyTorch} pipeline for \textbf{seven Hindustani raags} using the Saraga Hindustani dataset; implemented resampling, RMS normalisation and overlapping \textbf{30-second audio windows}.
  \item Explored per-window and full-recording tonic alignment against native-frequency \textbf{STFT/CNN} representations; trained a native-STFT CNN on \textbf{3,471 windows over 30 CUDA epochs}.
\end{cvitems}
\section{APPLIED EXPERIENCE}
\entry{Avartan Labs}{Product Engineering / Machine Learning Intern}{May--Jul 2026}
\begin{cvitems}
  \item Doubled rare-class correct detection from \textbf{40\% to 80\%} in an \textbf{18-class RT-DETR} sorting pipeline.
  \item Investigated focal loss, logit adjustment, class-specific thresholds and targeted fine-tuning for rare classes; evaluated detector failure modes through crop-based refinement.
\end{cvitems}
\section{PUBLICATION}
\textbf{Inference-Time Loss-Guided Colour Preservation in Diffusion Sampling} (2026).\\
Public preprint: \href{https://arxiv.org/abs/2601.17259}{arXiv:2601.17259}; colour-fidelity research described above.
\section{TECHNICAL SKILLS}
\textbf{Programming:} Python, C++, C, Bash, Lean.\quad \textbf{Frameworks:} PyTorch, Diffusers, TensorFlow.\\\textbf{Methods:} Optimal transport, diffusion models, reinforcement learning, computer vision, time-series analysis.
\end{document}
